% Einheit 7 von 8 – Term und Ereignis (bank/zufallsexperimente-und-pfadregeln/e7.jsonl, zusammenbau v0.1)
%% TODO Zeitmarke Einheit 7: Marken-Zeile nicht lesbar
%% TODO Zweigzeile Teil 1: was hier gelernt wird (Satz in Schülersprache aus dem Einheitstitel) – Einheit 7
\einheitenkopf[e7]{Einheit 7 von 8 · Term und Ereignis}
\zweigzeile{Abitur GK}
\setcounter{aufgabe}{26}

%% TODO Titel: Ich-kann-Satz fehlt in der Bank (Platzhalter: Kettenname) – Nr. 27
%% TODO Anweisung: ein Satz über der ersten Teilaufgabe fehlt in der Bank – Nr. 27
\begin{aufgabe}{Term und Ereignis}
%% TODO \rechenplatz{n}: Schrittzahl des Grundfalls fehlt in der Bank (nur bei fester Schreibform, 2.3 b)
\begin{teile}
\teil Im Term $4 \cdot 0{,}3 \cdot 0{,}7^3$: Wofür steht der Faktor 4?\\ \kreuz{eine Stufe}\\ \kreuz{eine Wiederholung}\\ \kreuz{die Zahl der Reihenfolgen}\\ \kreuz{das Gegenereignis}\\ \kreuz{mehrere Fälle}
\teil Ein Glücksrad zeigt Rot mit $\frac{1}{4}$, sonst Blau; es wird dreimal gedreht. Welches Ereignis hat die Wahrscheinlichkeit $\left(\frac{3}{4}\right)^3$?
\teil Ein Versuch gelingt mit $0{,}3$. Term und Wert für: erst zwei Fehlversuche, dann ein Erfolg. P = \leerfeld
\teil Ein Glücksrad liefert einen Gewinn mit $0{,}1$; es wird viermal gedreht. Welches Ereignis beschreibt $4 \cdot 0{,}1 \cdot 0{,}9^3$?
\teil Beschreibe ein Experiment und ein Ereignis zum Term $1 - 0{,}95^{10}$ und begründe.
\teil Eine Münze wird mehrmals geworfen. Ein Ereignis hat die Wahrscheinlichkeit $\frac{1}{32}$. Schreibe als Potenz und nenne ein passendes Ereignis.
\teil In einer Gruppe sind 12 Mädchen und 10 Jungen; fünf werden ausgelost. Welches Ereignis beschreibt „12 über 3“ · „10 über 2“ : „22 über 5“?
\teil Ein Versuch gelingt mit $0{,}35$; zwölf Versuche. P(genau fünf Erfolge) $=$ „12 über a“ $\cdot 0{,}35^b \cdot c^7$. Bestimme a, b und c. a = \leerfeld, b = \leerfeld, c = \leerfeld
\teil Würfel C trägt 2, 2, 2, 2, 4, 4; eine Münze zeigt 1 oder 6. Je Runde wird beides geworfen und addiert; vier Runden. Ereignis: alle vier Summen sind verschieden. Die Wahrscheinlichkeit ist $v \cdot \left(\frac{1}{2}\right)^w \cdot \left(\frac{1}{3}\right)^4$. Bestimme die natürlichen Zahlen v und w. \hfill (Abitur 2026 LK) \\ v = \leerfeld, w = \leerfeld
\end{teile}
\end{aufgabe}

%% TODO Titel: Ich-kann-Satz fehlt in der Bank (Platzhalter: Kettenname) – Nr. 28
%% TODO Anweisung: ein Satz über der ersten Teilaufgabe fehlt in der Bank – Nr. 28
\begin{aufgabe}{Term und Ereignis – Fehler finden, Begründen, Anwendung, Darstellungswechsel}
\begin{teile}
\teil Ein Glücksrad hat zehn gleiche Sektoren, drei gelb, sieben grün; es wird zweimal gedreht. Nina deutet $\left(\frac{7}{10}\right)^2$ als „zweimal Gelb“. Finde den Fehler.
\teil Begründe: Ein Produktterm ohne Vorfaktor beschreibt eine feste Reihenfolge.
\teil Ein Spamfilter erkennt eine Spam-Mail mit $0{,}98$. Eine Firma erhält an einem Tag 50 Spam-Mails. Was beschreibt $1 - 0{,}98^{50}$, und wie groß ist es?
\teil Ein Rad trifft mit $0{,}4$. Schreibe als Term und berechne: „Von drei Drehungen treffen genau zwei.“ P = \leerfeld
\end{teile}
\end{aufgabe}
